\documentclass{article}
\usepackage[utf8]{inputenc}
\usepackage{amsmath}
\usepackage{amsthm}
\usepackage{amssymb,amsfonts}
\usepackage{array}
\usepackage{mathrsfs}
\usepackage[hidelinks]{hyperref}
\usepackage[capitalize]{cleveref}
\usepackage{latexsym}
\usepackage{todonotes}
\usepackage{float}
\usepackage{nicefrac}
\usepackage{epsfig}
\usepackage{stmaryrd}
\usepackage{blkarray}
\usepackage{setspace}
\usepackage{tikz-cd}
\usepackage{xypic}
\usepackage{bbm,ifpdf,tikz}
\ifpdf
\usepackage{pdfsync}
\fi
\usetikzlibrary{positioning,matrix,arrows,decorations.pathmorphing}
\usepackage{enumitem}
\usepackage{extarrows}
\counterwithin{equation}{section}
\usepackage{leftindex}
\usepackage{subcaption}
\usepackage{adjustbox}


\usetikzlibrary{positioning,arrows}
\usetikzlibrary{decorations.markings}

\oddsidemargin=0pt
\evensidemargin=0pt
\topmargin=0in
\headheight=0pt
\headsep=0pt
\setlength{\textheight}{9in}
\setlength{\textwidth}{6.5in}

\makeatletter
\newtheorem*{rep@thm}{\rep@title}
\newcommand{\newrepthm}[2]{%
\newenvironment{rep#1}[1]{%
 \def\rep@title{#2 \ref{##1}}%
 \begin{rep@thm}}%
 {\end{rep@thm}}}
\makeatother
\newrepthm{theorem}{Theorem}
\newrepthm{proposition}{Proposition}
\newtheorem{theorem}{Theorem}[section]
\newtheorem{proposition}[theorem]{Proposition}
\newtheorem{claim}[theorem]{Claim}
\newtheorem{corollary}[theorem]{Corollary}
\newtheorem{lemma}[theorem]{Lemma}


\newtheorem{maintheorem}{Theorem}	
\renewcommand{\themaintheorem}{\Alph{maintheorem}}
\crefname{maintheorem}{theorem}{theorems}

\theoremstyle{definition}
\newtheorem{example}[theorem]{Example}
\newtheorem{definition}[theorem]{Definition}
\newtheorem{question}[theorem]{Question}
\newtheorem{conj}[theorem]{Conjecture}
\newtheorem{convention}[theorem]{Convention}
\newtheorem{remark}[theorem]{Remark}
\crefname{remark}{remark}{remarks}

\makeatletter
\newenvironment{subdefinition}[1]{%
  \def\subdefinitioncounter{#1}%
  \refstepcounter{#1}%
  \protected@edef\theparentnumber{\csname the#1\endcsname}%
  \setcounter{parentnumber}{\value{#1}}%
  \setcounter{#1}{0}%
  \expandafter\def\csname the#1\endcsname{\theparentnumber.\Alph{#1}}%
  \ignorespaces
}{%
  \setcounter{\subdefinitioncounter}{\value{parentnumber}}%
  \ignorespacesafterend
}
\makeatother
\newcounter{parentnumber}

\setenumerate{label=(\arabic*)}




\newcommand{\norm}[1]{\Vert #1 \Vert}
\newcommand{\vp}{\varphi}
\newcommand{\I}{\mathbf{i}}
\newcommand{\J}{\mathbf{j}}

\newcommand{\bin}[2]{
\left(
 \begin{array}{@{}c@{}}
    #1 \\ #2
\end{array}
\right)  }
\newcommand{\td}[1]{\widetilde{#1}}
\newcommand{\ol}[1]{\overline{#1}}
\newcommand{\ora}[1]{\overrightarrow{#1}}
\newcommand{\into}{\hookrightarrow}
\newcommand{\onto}{\twoheadrightarrow}

\newcommand{\myeq}[1]{\mathrel{\overset{\makebox[0pt]{\mbox{\normalfont\tiny\sffamily #1}}}{=}}}



\newcommand{\Rbar}{\overline{\mathbb{R}}}
\newcommand{\C}{\mathbb{C}}
\newcommand{\bM}{\mathbb{M}}



\newcommand{\rk}{\operatorname{rk}}
\newcommand{\cl}{\operatorname{cl}}
\newcommand{\rL}{\mathrm{L}}
\newcommand{\rT}{\mathrm{T}}
\newcommand{\epn}{\mathtt{exp}}
\newcommand{\tM}{\mathtt{M}}
\newcommand{\ds}{\oplus}
\newcommand{\Tr}{\mathrm{Tr}}
\newcommand{\stcap}{\cap_{st}}
\newcommand{\Trop}{\operatorname{Trop}}
\newcommand{\trop}{\operatorname{trop}}
\newcommand{\val}{\operatorname{val}}
\newcommand{\vol}{\operatorname{vol}}
\newcommand{\Vol}{\operatorname{Vol}}
\newcommand{\nk}[2]{\binom{#1}{#2}}
\newcommand{\supp}{\operatorname{supp}}
\newcommand{\PG}{\operatorname{PG}}


\newcommand{\jayden}[1]{{ \sf $\clubsuit\clubsuit\clubsuit$ {\textcolor{red}{Jayden: [#1]}}}}

\newcommand{\sam}[1]{{ \sf $\diamondsuit\diamondsuit\diamondsuit$ {\textcolor{red}{Sam: [#1]}}}}

\newcommand{\rec}{\operatorname{rec}}
\newcommand{\Star}{\operatorname{Star}}
\newcommand{\op}{\operatorname{op}}
\newcommand{\sed}{\operatorname{sed}}
\newcommand{\cone}{{\operatorname{cone}}}
\newcommand{\sym}{{\operatorname{sym}}}

\renewcommand{\deg}{\operatorname{deg}}
\newcommand{\cp}{\mathbb{C}P}


\newcommand{\Z}{\mathbb{Z}}
\newcommand{\Q}{\mathbb{Q}}
\newcommand{\N}{\mathbb{N}}
\newcommand{\R}{\mathbb{R}}
\newcommand{\bF}{\mathbb{F}}
\newcommand{\K}{\mathbb{K}}
\newcommand{\T}{\mathbb{T}}
\newcommand{\PP}{\mathbb{P}}
\newcommand{\PT}[1]{\mathbf{P}\big(\mathbb{\Rbar}^{#1}\big)}

\newcommand{\cE}{\mathcal{E}}
\newcommand{\cO}{\mathcal{O}}
\newcommand{\cW}{\mathcal{W}}
\newcommand{\cV}{\mathcal{V}}
\newcommand{\cX}{\mathcal{X}}
\newcommand{\cN}{\mathcal{N}}
\newcommand{\cHH}{\mathcal{H}}
\renewcommand{\cR}{\mathcal{R}}
\newcommand{\cP}{\mathcal{P}}
\newcommand{\cG}{\mathcal{G}}
\newcommand{\cF}{\mathcal{F}}
\newcommand{\cC}{\mathcal{C}}
\newcommand{\cT}{\mathcal{T}}
\newcommand{\cJ}{\mathcal{J}}
\newcommand{\cK}{\mathcal{K}}
\newcommand{\cI}{\mathcal{I}}
\newcommand{\cA}{\mathcal{A}}
\newcommand{\cS}{\mathcal{S}}
\newcommand{\cM}{\mathcal{M}}
\newcommand{\cQ}{\mathcal{Q}}
\newcommand{\cB}{\mathcal{B}}
\renewcommand{\cD}{\mathcal{D}r}
\newcommand{\cZ}{\mathcal{Z}}
\newcommand{\cLL}{\mathcal{L}}

\newcommand{\sP}{\mathscr{P}}
\newcommand{\sF}{\mathscr{F}}

\newcommand{\tB}{\mathtt{B}}
\newcommand{\tD}{\mathtt{D}}

\newcommand{\Gr}{\mathbf{Gr}}
\newcommand{\Dr}{\mathbf{Dr}}
\newcommand{\Qt}{\mathbf{Qt}}
\newcommand{\bQ}{\mathbf{Q}}
\newcommand{\bC}{\mathbf{C}}
\newcommand{\be}{\mathbf e}
\newcommand{\bv}{\mathbf v}
\newcommand{\bV}{\mathbf V}
\newcommand{\bu}{\mathbf u}
\newcommand{\ba}{\mathbf a}
\newcommand{\bb}{\mathbf b}
\newcommand{\bc}{\mathbf c}
\newcommand{\bx}{\mathbf x}
\newcommand{\by}{\mathbf y}
\newcommand{\bw}{\mathbf w}
\newcommand{\bW}{\mathbf W}
\newcommand{\bh}{\mathbf h}
\newcommand{\br}{\mathbf r}

\newcommand{\rmL}{\mathrm{L}}
\newcommand{\rmS}{\mathrm{S}}
\newcommand{\DT}{\operatorname{\mathrm{DT}}}
\newcommand{\rmAd}{\mathrm{Ad}}
\newcommand{\rmH}{\mathrm{H}}
\newcommand{\rmV}{\mathrm{V}}


\author{Sam Payne \and Jidong Wang}

\title{First proof: Duality for tropical linear subspaces}
\date{}

\begin{document}

\maketitle


Let $M$ be a matroid, and let $\Trop(M)$ be the corresponding tropical linear space, i.e., the closure of the Bergman fan of $M$ in tropical projective space. The relative Dressian of $M$ is 
\[ \Dr(M) :=\{\text{ tropical linear subspaces of $\Trop M$ }\}, \]
which is partially ordered by inclusion of tropical linear spaces. 

Suppose $M$ admits an order-reversing involution $F\mapsto F^\perp$ on the lattice of flats $\cLL(M)$. Note that there is a natural inclusion $\cLL(M)^{\op} \subset \Dr(M)$ given by $F\mapsto \Trop(M/F\oplus U_{0,F})$.

\medskip
\noindent \textbf{Question.}  Is there an order-reversing involution $\Phi$ on $\Dr(M)$ that extends the involution $F\mapsto F^\perp$?%

\medskip

\noindent \textbf{Answer.} Yes. There is such an order-reversing involution. Further details will appear in \cite{wang26modular}.


\section{Context and motivation}



Tropical linear subspaces of the tropical projective space $\PT{n}$ correspond naturally and bijectively to valuated matroids on $[n]$, up to tropical rescaling, and a valuated matroid of rank $k$ on $[n]$ has a valuated dual of rank $n-k$. The resulting involution is order-reversing on tropical linear subspaces of $\PT{n}$.

One motivating example is the Boolean matroid $U_{n,n}$. Its lattice of flats $\cLL(U_{n,n})$ is the lattice of subsets of $[n]$, and $\Trop U_{n,n}$ is the tropical projective space $\PT{n}$. Duality for valuated matroids induces an involution on $\Dr(U_{n,n})$ that extends the natural involution $S \mapsto S^c$ on $\cLL(U_{n,n})$.



Another motivation comes from the classical geometry of linear subspaces and Grassmannians. Let $V$ be a finite-dimensional vector space. A nondegenerate pairing on $V$ induces an isomorphism $V\xrightarrow{\sim}V^\vee$, sending $v$ to $v^\perp$. The map $W\mapsto W^\perp$ is an order-reversing involution on $\bigsqcup_k\Gr(k,V)$. 

We note that matroids whose lattices of flats admit an order-reversing involution are called \emph{modular matroids}. One does not expect any natural duality for tropical linear subspaces of tropical linear spaces that correspond to non-modular matroids. 




\medskip

\section{Proof strategy}
A valuated matroid $\mu$ of rank $k$ on $[n]$ is a function $\nk{n}{k}\to \Rbar$ that satisfies the tropical Plücker relations:
\begin{itemize}
    \item for each $S\subset [n]$ of size $k-1$ and each $T\subset [n]$ of size $k+1$, 
\begin{equation}\label{eqn:plucker}
    \min_{a\in T\smallsetminus S}\{\mu(T-a )+\theta(S+a)\}
\end{equation}
vanishes tropically, i.e., either all terms in \eqref{eqn:plucker} are infinity, or the minimum is achieved at least twice.
\end{itemize}
Here, $T-a:=T\smallsetminus\{a\}$ and $S+a:=S\cup\{a\}$. We call \eqref{eqn:plucker} the \emph{Pl\"ucker relation labeled by $(S,T)$}.

Tropical linear spaces correspond naturally and bijectively to valuated matroids, up to tropical scaling. For each valuated matroid $\theta$ of rank $k$, let $\Trop\theta$ be the corresponding tropical linear space of dimension $k-1$. A valuated matroid $\theta_1$ is a quotient of $\theta_2$, denoted $\theta_2\onto \theta_1$, if they satisfy the tropical incidence-Plücker relations \eqref{eqn:incidence-plucker}, and $\theta_2\onto\theta_1$ if and only if $\Trop\theta_2\supset \Trop\theta_1$. We define 
\[\bV^k(\mu):=\{\text{ corank $k$ valuated matroid quotients of $\mu$ }\}.\]
By construction, $\bV^k(\mu)\subset \Rbar^{\nk{n}{r-k}}$, where $r$ is the rank of $\mu$. 

Every matroid $M$ gives rise to a valuated matroid, called the \textit{trivial valuation on $M$}, that sends $B\in \nk{n}{r}$ to 0 if $B$ is a basis of $M$, and to $\infty$ otherwise. We denote by $\bV^k(M)$ the set of valuated matroid quotients of the trivial valuation on $M$ and let $\bV(M)$ be the disjoint union $\bigsqcup_k\bV^k(M)$. Note that $\bV(M)$ is partially ordered by the valuated matroid quotient relation.  We construct a map $\theta\mapsto\theta'$ on $\bV(M)$ that induces an order-reversing involution on $\Dr(M)$. The map is relatively simple, and can be viewed as the restriction of a tropical linear map in the natural Pl\"ucker coordinates.

In the analogous classical setting, the Pl\"ucker coordinates of $[W]\in\Gr(k,\bF^r)$ are in bijection with coordinate subspaces of dimension $k$ in $\bF^r$. The matroid analog of a dimension $k$ coordinate subspace is a rank $k$ flat. Therefore, one might naturally guess that there should be a cryptomorphic characterization of valuated matroid quotients of rank $k$ as functions on the rank $k$ flats of $M$. We construct such a cryptomorphism in \Cref{prop:essential-relations}. The map $\Phi$ is then given simply by $\theta\mapsto\theta'$, where $\theta'(F):=\theta(F^\perp)$.

\section{Solution}

Throughout, $M$ denotes a simple matroid on $[n]$. 

\subsection{A cryptomorphism of matroid quotients}

Let $N$ be a matroid on $[n]$ with rank function $\rk_N$. Then $N$ is a matroid quotient of $M$ if and only if 
\begin{equation}\label{eqn:quotient-rank-function}
    \rk_N(A+x)-\rk_N(A)\leq \rk_M(A+x)-\rk_M(A)
\end{equation}
for any $A\subset[n]$ and any $x\in [n]$. \Cref{lem:quotient-rank-function} is a cryptomorphic characterization of matroid quotients in terms of rank functions on $\cLL(M)$. The proof is straightforward and is omitted. 
\begin{lemma}\label{lem:quotient-rank-function}
    Let $\rho \colon \cLL(M)\to \N$ be any function that satisfies the following three properties
    \begin{itemize}
    \itemsep0em
        \item \emph{(Submodularity)} \hspace{.2 em}  $\rho(F)+\rho(G)\geq \rho(F\vee G)+\rho(F\cap G)$;
        \item \emph{(Unit increment)} \hspace{.1 em} $\rho(F\vee x)-\rho(F)\in \{0,1\}$;
        \item \emph{(Normalization)}  \hspace{.4 em} $\rho(\emptyset)=0$.    \end{itemize}
        Then $A\mapsto\rho(\cl_M(A))$ is the rank function of a matroid quotient of $M$. Conversely, if $\rho$ is the rank function of a matroid quotient of $M$, then the restriction of $\rho$ to $\cLL(M)$ satisfies these three properties.
\end{lemma}
\noindent The proof is straightforward and omitted. We call the functions $\cLL(M)\to\N$ satisfying the properties in \Cref{lem:quotient-rank-function} \emph{lattice rank functions}. 

Let $\bQ^k(M)$ be the set of corank $k$ matroid quotients of $M$, with  $\bQ(M) := \bigsqcup_k\bQ^k(M)$, partially ordered by the matroid quotient relation. 

\begin{proposition}\label{prop:existence-Qt-dual}
Let $M$ be a modular matroid and $\rho\colon\cLL(M)\to \N$ be its lattice rank function. Then for any $N\in \bQ^k(M)$ with lattice rank function $\rk_N$, the function $\cLL(M)\to \N$ given by
\begin{equation}\label{eqn:dual-rank-function}
    F \quad \mapsto\quad  \rho(F)-\rk_N([n])+\rk_N(F^\perp),
\end{equation}
is the lattice rank function of some $N'\in \bQ^{r-k}(M)$. The resulting map $N\mapsto N'$ is an order-reversing involution on $\bQ(M)$, and $\Psi(M/F\oplus U_{0,F})=M/F^\perp \oplus U_{0,F^\perp}$. 
\end{proposition}

\begin{proof}
    The proof that \eqref{eqn:dual-rank-function} is the lattice rank function of some corank $k$ quotient $N'\in \bQ^k(M)$ is identical to the proof of \cite[Theorem~42]{MR3829288}. A straightforward computation shows $(N')'=N$. Next, we show that $N'_2\onto N'_1$ whenever $N_1\onto N_2$. For any $F\subset G\in \cLL(M)$, we compute the difference 
    \begin{align*}
        \rk_{N'_2}(F)-\rk_{N'_2}(G) - (\rk_{N'_1}(F)-\rk_{N'_1}(G)) = \rk_{N_1}(F^\perp) - \rk_{N_1}(G^\perp) - (\rk_{N_2}(F^\perp)-\rk_{N_2}(G^\perp)).
    \end{align*}
    This is nonnegative because $N_1\onto N_2$, and hence $N'_2\onto N'_1$ by \eqref{eqn:quotient-rank-function}.

    If $\rk$ is the lattice rank function of $M/F\oplus U_{0,F}$, then we have
    \[\rk(G)=\rho(G\vee F)-\rho(F).\]
    Then \eqref{eqn:dual-rank-function} maps $G$ to $\rho(G)+\rk(G^\perp)-(n-\rho(F))$. Collecting terms, $G\mapsto \rho(G\vee F^\perp)-\rho(F^\perp)$, which is the lattice rank function of $M/F^\perp\oplus U_{0,F^\perp}$.
\end{proof}

The following observation is immediate from the construction of $\Psi$.
\begin{proposition}\label{prop:basis-M-dual}
   The bases of $N'$ are $\{B\mid \cl_M(B)^\perp=\cl_M(D)\text{ for some basis $D$ of $N$}\}$.
\end{proposition}

\subsection{A cryptomorphism of valuated matroid quotients}

 Let $\mu$ be a valuated matroid of rank $r$ on $[n]$, whose underlying matroid is $M$. A valuated matroid $\theta$ of rank $s$ on $[n]$ is a quotient of $\mu$ if and only if $\mu$ and $\theta$ satisfies the following tropical incidence-Pl\"ucker relations: for any $S\subset [n]$ of size $s-1$ and $T\subset [n]$ of size $r+1$,
\begin{equation}\label{eqn:incidence-plucker}
    \min_{a\in S\smallsetminus T}\{\mu(T-a)+\theta(S+a)\}
\end{equation}
 vanishes tropically, i.e., either every term in \eqref{eqn:incidence-plucker} is infinity, or the minimum is achieved at least twice. We say \eqref{eqn:incidence-plucker} is the \textit{incidence relation} \textit{labeled by} $(S,T)$.


In what follows, we adopt a common convention from the literature on matroids and Grassmannians and use juxtapositions to denote unions, e.g. $SDab := S \cup D \cup \{a,b\}$. We will also need the following equivalent formulation of the tropical Plücker relations \cite{brandt2021tropical}, called the \textit{valuated basis exchange}. For any pair of bases $(B_1,B_2)$ of $\mu$, and for any $a\in B_1\smallsetminus B_2$, there is some $b\in B_2\smallsetminus B_1$ such that 
    \[\mu(B_1)+\mu(B_2)\geq \mu(B_1-a+b)+\mu(B_2-b+a).\]

\begin{proposition}\label{prop:essential-relations}
Let $\mu$ be a valuated matroid on $[n]$ of rank $r$ with underlying matroid $M$ and $\theta$ be a valuated matroid on $[n]$ of rank $r-k$ with underlying matroid $N$. Then $\theta\in \bV^k(\mu)$ if and only if it has the following properties. 

     \begin{itemize}
     \itemsep0em
         \item \emph{(Degenerate flats)} \hspace{0 em} If $\rk_M(A)<r-k$, then $\theta(A)=\infty$.
         \item \emph{(Parallel flats)} \hspace{1.5 em} If  $\cl_M(A_1)=\cl_M(A_2)$, and $A_1D$ and $A_2D$ are bases of $M$, then $$\mu(A_1 D)-\mu(A_2 D)=\theta(A_1)-\theta(A_2).$$ 
         \item \emph{(Concurrent flats)} \hspace{0 em} Let $S \subset [n]$ be a subset of size $r-k-1$ that is  independent in $M$. For each circuit $C$ of each contraction $M/S$, let $D_C$ be a spanning set of $M/S$ of size $k+2$ containing $C$. Then
         \begin{equation}\label{eqn:concurrent-flats}
             \min_{a\in C}\{\theta(S a)+\mu(S D_C -a)\}
         \end{equation}
         vanishes tropically. 
     \end{itemize}
 \end{proposition}
 
     \begin{proof}
         Suppose $\theta\in \bV^k(\mu)$. We show $\theta$ satisfies the three properties in the proposition. 

        \begin{itemize}
            \item (Degenerate flats)  Let $C_0$ be a circuit in $A$. Let $D_0$ be a spanning set of $M$ containing $C_0$ of size $r+1$. Pick $x\in C_0$. The terms in the incidence relation labeled by $(A-x,D_0)$ are all infinity except for possibly $\theta(A)+\mu(D-x)$ so $\theta(A)=\infty$.

            \item (Parallel flats)  Using basis exchange, it suffices to prove the case $|A_1\smallsetminus A_2|=1$. Assume $A_1=Sa$ and $A_2=Sb$. Let $D_1$ be such that $D_1 Sa$ is a basis of $M$. Then the parallel flats property follows from the incidence relation \eqref{eqn:incidence-plucker} labeled by $(S, S Dab)$.

            \item (Concurrent flats)  Let $S\in \nk{n}{r-k-1}$ be independent. Let $C$ be a circuit of the of contraction $M/S$ and $D$ be a spanning set of $M/S$ of size $k+2$ containing $C$. Then \eqref{eqn:concurrent-flats} is simply the incidence relation labeled by $(S,S D)$. 
            \end{itemize}
         
         
        
         
         Now, suppose $\theta$ is a valuated matroid that satisfies the three properties. Note that the parallel flats property and the concurrent flats property are exactly the nontrivial incidence relations labeled by $(S,S T)$ for some $S$ independent in $M$ and some $T$. 
         
         Consider an arbitrary incidence relation for $\mu\onto\theta$ labeled by $(S,D)$. We may assume $S$ is independent in $M$, and that $D$ does not contain a circuit of $M$ that is contained in $S$. For contradiction, suppose the relation is not satisfied and the unique minimum is achieved uniquely by $\theta(S1)+\mu(D-1)<\infty$. Since $S1$ is independent and $D-1$ is a basis of $M$, there is some $T_1\subset D-1$ such that $S1T_1$ is a basis of $M$. Apply valuated basis exchange for $\mu$ to the pair of bases $(S1T_1,D-1)$, we get some $c_1\in D-1$ such that 
    \[\mu(S1T_1)+\mu(D-1)\geq \mu(ST_1c_1)+\mu(D-c_1).\]
    All the values above are finite, so we can write 
    \[\mu(S1T_1)-\mu(ST_1c_1)\geq \mu(D-c_1)-\mu(D-1).\]
    Adding $\theta(Sc_1)-\theta(S1)$ (which, a priori, may be infinite) to both sides, we get
    \[\theta(Sc_1)+\mu(S1T_1)-\theta(S1)-\mu(ST_1c_1)\geq \theta(Sc_1)+\mu(D-c_1)-\theta(S1)-\mu(D-1)>0,\]
    so 
    \[\theta(Sc_1)+\mu(S1T_1)>\theta(S1)+\mu(ST_1c_1).\]
    This implies, by the incidence relation labeled by $(S,S1T_1c_1)$, that there is some $b_1\in T_1$ such that
    \[\theta(Sc_1)+\mu(S1T_1)>\theta(S1)+\mu(ST_1c_1)\geq \theta(Sb_1)+\mu(S1T_1+c_1-b_1).\]
    Set $T_2=T_1+c_1-b_1$. Since $\theta(Sb_1)+\mu(ST_1c_1)<\infty$, $\theta(Sb_1)+\mu(S1T_2)<\infty$. Suppose we have obtained $T_\ell$ this way. Since $\mu(S1T_\ell)<\infty$, we can apply valuated basis exchange to the pair $(S1T_\ell, D-1)$, we get some $c_\ell\in D-1$ such that 
    \[\mu(S1T_\ell)+\mu(D-1) \geq \mu(ST_\ell c_\ell) + \mu(D-c_\ell),\]
    where all the terms are finite. Again, this implies 
    \[\theta(Sc_\ell)+\mu(S1T_\ell)>\theta(S1)+\mu(ST_\ell c_\ell)\]
    and by the incidence relation labeled by $(S,S1T_\ell c_\ell)$, there is some $b_\ell\in T_\ell$ such that 
    \[\theta(Sc_\ell)+\mu(S1T_\ell)>\theta(S1)+\mu(ST_\ell c_\ell)\geq \theta(Sb_\ell)+\mu(S1T_\ell+c_\ell-b_\ell).\]
    Now set $T_{\ell+1}=T_\ell+c_\ell-b_\ell$. For each $\ell$, we have inequality
    \[\theta(Sc_\ell)+\mu(S1T_\ell)>\theta(Sb_\ell)+\mu(S1T_{\ell+1}),\]
    where the right hand side is finite. Since $T_\ell\subset D$ for any $\ell$, there are only finitely many possibilities for $T_\ell$, so for some $\ell_2>\ell_1$, we have $T_{\ell_1}=T_{\ell_2}$. We have inequality of the sum
    \begin{equation}\label{eqn:inequality-contradiction}
        \sum_{i=\ell_1}^{\ell_2-1}\Big(\theta(Sc_i) + \mu(S1T_i)\Big) > \sum_{i=\ell_1}^{\ell_2-1}\Big(\theta(Sb_i) + \mu(S1T_{i+1})\Big). 
    \end{equation}
    Since $T_{i+1} b_i = T_i c_i$, we have equality of multi-sets
    \[\bigsqcup_{i=\ell_1}^{\ell_2-1}(T_{i+1} b_i) = \bigsqcup_{i=\ell_1}^{\ell_2-1} (T_i c_i).\]
    Thus $\{b_{\ell_1},...,b_{\ell_2-1}\}= \{c_{\ell_1}, ..., c_{\ell_2-1}\}$ as multi-sets, and the two sides of \eqref{eqn:inequality-contradiction} are equal, a contradiction.
\end{proof}



     



\subsection{Construction of \texorpdfstring{$\Phi$}{Phi}}
 
For each $\theta\in \bV^k(M)$, we define a function $\Phi(\theta) \colon \cLL^{r-k}(M) \to \Rbar$ by 
\[I\mapsto \begin{cases}
    \theta(J), & \text{ if }\cl_M(J)=\cl_M(I)^\perp,\\
    \infty, & \text{ otherwise}.
\end{cases}\]
It is clear that $\Phi$ is compatible with $F \mapsto F^\perp$.  We now show that $\Phi(\theta)$ is a rank $k$ valuated quotient of $M$.



\noindent 

The tropical linear space $\Trop M$ is the intersection of the tropical hyperplanes corresponding to the circuits of $M$. A \textit{tropical basis} of $M$ is a subset of the circuits whose corresponding tropical hyperplanes cut out $\Trop M$. Recall the \textit{pasting} operation from \cite{yu2006representing}: if $C_1$ and $C_2$ are two circuits of $M$ with $C_1\cap C_2=\{a\}$ then pasting them means taking their symmetric difference $(C_1\cup C_2)\smallsetminus\{a\}$. If every circuit in $M$ can be constructed by successive pasting of circuits in a collection $\cS$, then $\cS$ is a tropical basis \cite[Lemma 4]{yu2006representing}. 

\begin{proposition}\label{prop:modular-tropical-basis}
    For any simple modular matroid, the 3-element circuits form a tropical basis.
\end{proposition}

\begin{proof}
    Every simple modular matroid is a direct sum of a free matroid and finite projective geometries. We only need to prove the statement for $M$ a connected finite projective geometry. If $M=U_{2,n}$, all circuits have size 3. Assume $M$ is a finite projective geometry of rank $\geq 3$. We will 
    show that the 3-element circuits of $M$ generate all circuits via pasting. We proceed by induction on the rank of $M$.

    If $M$ has rank 3, then a size 4 circuit $\{a_1,a_2,b_1,b_2\}$ is a symmetric difference between two circuits $\{a_1,a_2,c\}$ and $\{b_1,b_2,c\}$, where $c$ is the point where the line $\cl_M(a_1a_2)$ and the line $\cl_M(b_1b_2)$ intersect. Let $M$ have rank $r>3$. By the induction hypothesis, any circuit of size $\leq r$, which is a circuit of a finite projective geometry of lower rank, can be obtained by pasting the 3-element circuits. Let $C$ be a circuit of size $r+1$. Take two distinct points $a,b\in C$ and let $C'=C\smallsetminus\{a,b\}$. The line $\cl_M(ab)$ intersects with $\cl_M(C')$ at a point $d$. Note that $d$ is not in the closure of any proper subsets of $C'$, otherwise $C$ could not be a circuit. This means both $C'\cup \{d\}$ and $\{a,b,d\}$ are  circuits. Since $C$ is the symmetric difference of $C'$ and $\{a,b,d\}$, the claim follows by induction.
\end{proof}

\begin{proposition}\label{prop:3-term-incidence}
    If $\mu$ is the trivial valuation on a modular matroid $M$, then the concurrent flats property in \Cref{prop:essential-relations} can be replaced by the following condition,
    \begin{itemize}
        \item \emph{(Concurrent triples)} Let $S\in \nk{n}{r-k-1}$ be independent in $M$. For each circuit $C=\{a,b,c\}$ of $M/S$, 
        \begin{equation}\label{eqn:concurrent-triple-relation}
            \min\{\;\theta(Sa),\quad \theta(Sb), \quad \theta(Sc)\;\}
        \end{equation}
        vanishes tropically.
    \end{itemize}
    In other words, a valuated matroid $\theta$ on $[n]$ of rank $r-k$ is in $\bV(M)$ if and only if it satisfies the degenerate flats property, the parallel flats property, and the concurrent triples property.
\end{proposition}

\begin{proof}
The contraction $M/S$ is a modular matroid, so the statement follows from \Cref{prop:modular-tropical-basis}.
\end{proof}

We say that three distinct flats $F_1,F_2,F_3\in \cLL(M)$ form a \textit{concurrent triple} if they cover the same flat $G$. The following corollary follows easily from the parallel flats property and the concurrent triples property.

\begin{corollary}\label{cor:concurrent-triple}
    Let $A_1,A_2$ and $A_3$ be independent sets of $M$ such that $\cl_M(A_1),\cl_M(A_2)$ and $\cl_M(A_3)$ form a concurrent triple of corank $k$ flats. Then for any $\theta\in \bV^k(M)$,
    \begin{equation}
            \min\{\;\theta(A_1),\quad \theta(A_2), \quad \theta(A_3)\;\}
        \end{equation}
        vanishes tropically.
\end{corollary}


Let $N$ be the underlying matroid of $\theta \in \bV^k(M)$. Denote $\Phi(\theta)$ by $\theta'$. Note that the support of $\theta'$ is the set of bases of $N'$, as in Proposition~\ref{prop:existence-Qt-dual}. 

The degenerate flats properties and the parallel flats properties for $\theta'$ are clear. To show that $\theta'$ is in $\bV^{r-k}(M)$, it is enough to verify that it satisfies  concurrent triples properties, by Proposition~\ref{prop:3-term-incidence}.   Suppose $A\in \nk{n}{k-1}$ is independent in $M$ and $C=\{1,2,3\}$ is a 3-element circuit of the contraction $M/A$. Let $F:=\cl_M(A)$, $F_i:=\cl_M(Ai)$ for each $i\in\{1,2,3\}$, and $F_{123}=\cl_M(A123)$. Then $F_1,F_2$ and $F_3$ form a concurrent triple. By the modularity of $M$, $F_1^\perp$, $F_2^\perp$,  and $F_3^\perp$ form a concurrent triple. It then follows from \Cref{cor:concurrent-triple} that
        \[\min\{\;\theta' (A_1),\quad \theta' (A_2),\quad \theta' (A_3)\}\]
        vanishes tropically.
        
        Now, we show that $\theta'$ is a valuated matroid. Since the support of $\theta'$ is a matroid, it suffices to verify that $\theta'$ satisfies all the 3-term Plücker relations \cite[Theorem 5.2.25]{murota2010matrices}. Choose $S\in \nk{n}{k-2}$ and, by relabeling, assume $1,2,3,4\notin S$. We need to show
    \begin{equation}\label{eqn:plucker-to-prove}
        \min\{\theta'(S12)+\theta'(S34),\theta'(S13)+\theta'(S24),\theta'(S14)+\theta'(S23)\}
    \end{equation}
    vanishes tropically. Define
     \begin{equation}\label{eqn:flats}
         H:=\cl_M(S),\quad H_a:=\cl_M(Sa),\quad H_{ab}:=\cl_M(Sab),\quad H_{1234}:=\cl_M(S1234),
     \end{equation}
     where $a,b,c,d\in \{1,2,3,4\}$ are distinct elements. Up to relabeling, we may assume that $S12$ and $S34$ are both independent. We carried out the cases according to the contraction $M'=M|_{S1234}/S$. Let $\cP$ be the set of all the flats in \eqref{eqn:flats}. 
    \begin{itemize}
        

        \item Case 2: $M'\cong U_{3,4}$. Then any $\{b,c,d\}\subset \{1,2,3,4\}$ is a circuit of $M|_{Sabcd}/Sa$. We have just shown that 
        \[\min\{\theta'(Sab),\quad \theta'(Sac),\quad  \theta'(Sad)\}\]
        vanishes tropically by the concurrent triple property for any distinct $a,b,c,d\in \{1,2,3,4\}$. Now it follows from \cite[Theorem 3.10]{joswig2023generalized} that \eqref{eqn:plucker-to-prove} vanishes tropically.

        \item Case 3: $M'\cong U_{2,2}\oplus U_{1,2}$. We may assume $H_{12}=H_1=H_2$. Then $\theta'(S12)=\infty$ and \eqref{eqn:plucker-to-prove} becomes 
        \[\min\{\theta'(S13)+\theta'(S24),\;\theta'(S14)+\theta'(S23)\},\]
        which vanishes tropically, because $\theta'(S13)=\theta'(S23)$ and $\theta'(S14)=\theta'(S24)$.

        \item Case 4: $M'\cong U_{2,3}\oplus U_{1,1}$. We may assume $H_{12}=H_{13}=H_{23}$. Then $\theta'(S12)=\theta'(S13)=\theta'(S23)$. Note that  \eqref{eqn:plucker-to-prove} follows, because the concurrent triples property says
        \[\min\{\theta'(S14),\quad \theta'(S24),\quad, \theta'(S14)\}\]
        vanishes tropically.
        
        \item Case 5: $M'$ has rank 2. Recall that we assume $S12$ and $S34$ are both independent. By the basis exchange property, we may also assume, by relabeling, that $S13$ and $S24$ are both independent. Then we have  $H_{12}=H_{34}=H_{13}=H_{24}$. Either $\cl_M(S14)=H_{12}$ or $S14$ is dependent. Hence, $\theta'(S14)$ is either equal to $\theta'(S12)$ or infinity, meaning \eqref{eqn:plucker-to-prove} vanishes tropically.
    \end{itemize}


It remains to show that $\Psi$ reverses the partial order on $\bV(M)$.



\begin{proposition}\label{prop:order-reversing}
         For any $\theta_1\in \bV_t(M)$ and $\theta_2\in \bV_s(M)$ where $t\geq s$, $\theta_1\onto \theta_2$ if and only if $\Phi(\theta_2) \onto \Phi(\theta_1)$.
\end{proposition}

\begin{proof}
Denote $\Phi(\theta)$ by $\theta'$. Since $\Phi$ is an involution on $\bV(M)$, it suffices to show that $\theta_2'\onto \theta_1'$ implies $\theta_1\onto \theta_2$. By \Cref{prop:existence-Qt-dual} and \Cref{prop:basis-M-dual}, the underlying matroid of $\theta_2$ is a quotient of the underlying matroid of $\theta_1$. This means the degenerate flats property is satisfied. By \Cref{prop:essential-relations}, it remains to show that the parallel flats property and the concurrent flats property are satisfied. Recall that those are exactly the incidence relations for $\theta_1\onto \theta_2$ labeled by $(S,S T)$ where $|S|=s-1$ and $|T|=t-s+2$. By \Cref{prop:essential-relations}, we may assume $\rk_M(S T)\geq t$ for the relation to be nontrivial.
\begin{itemize}
    \item Case 1: $\rk_M(S T)=t$. There is a unique circuit $C$ of $M|_{ST}/S$ and $T$ is a spanning set of $M|_{ST}/S$ that contains $C$. Since $\cl_M(S T-a)=\cl_M(S T)$ for any $a\in C$, by the parallel flats property, $\theta_1(ST-a)$ does not depend on $a\in C$. Hence
    \[\min_{a\in C}\{\theta_1(S T-a)+\theta_2(S+a)\}\]
    vanishes tropically, by the concurrent flats property satisfied by $\theta_2$ that $\min_{a\in C}\{\theta_2(S+a)\}$ vanishes tropically.

    \item Case 2: $\rk_M(S T)=t+1$. For each $I\subset T$, let $F_I=\cl_M(S I)$. Notice that $S T$ is a basis of $\rk_M(S T)$, so $F_I\cap F_J=F_{I\cap J}$ and one has a Boolean sublattice
    \[\cP:=\{F_I\mid I\subset T\}.\]
    The dual of this sublattice 
    \[\cP^\perp:=\{F_I^\perp \mid I\subset T\}\]
    is also a Boolean sublattice. Let $S'$ be a basis of $F_\emptyset^\perp$. Then there is some $T'$ of size $t-s+2$ such that $\cP^\perp$ is the Boolean sublattice 
    \[\{G_I:=\cl_M(S' I)\mid I\subset T'\}.\]
    We thus obtain a bijection $\sigma \colon T \to T'$ by setting $F_{\{a\}}\mapsto F_{\{a\}}^\perp = G_{T'\smallsetminus\{\sigma(a)\}}$. By definition, $\theta_1(ST-a)=\theta_1'(S'+\sigma(a))$ and $\theta_2(S+a)=\theta_1'(S'T'-\sigma(a))$. Hence, the incidence relation for $\theta_1\onto\theta_2$ labeled by $(S,ST)$ is the same as the incidence relation labeled for $\theta_2'\onto \theta_1'$ by $(S',S'T')$, and is thus satisfied.
\end{itemize}
Since in either case, the incidence relations for $\theta_1\onto \theta_2$ are satisfied, the statement is proved.
\end{proof}

We have shown that $\Phi$ maps $\bV^k(M)$ to $\bV^{r-k}(M)$ and is order-reversing. It is then straightforward to check that $\Phi$ is also an involution, which completes the solution.


\providecommand{\bysame}{\leavevmode\hbox to3em{\hrulefill}\thinspace}
\providecommand{\MR}{\relax\ifhmode\unskip\space\fi MR }

\providecommand{\MRhref}[2]{%
  \href{http://www.ams.org/mathscinet-getitem?mr=#1}{#2}
}
\providecommand{\href}[2]{#2}
\begin{thebibliography}{JLLAO23}

\bibitem[BEZ21]{brandt2021tropical}
Madeline Brandt, Christopher Eur, and Leon Zhang, \emph{Tropical flag
  varieties}, Adv. Math. \textbf{384} (2021), 107695.

\bibitem[FM19]{Fink_2019}
Alex Fink and Luca Moci, \emph{Polyhedra and parameter spaces for matroids over
  valuation rings}, Adv. Math. \textbf{343} (2019), 448--494.

\bibitem[JLLO23]{joswig2023generalized}
Michael Joswig, Georg Loho, Dante Luber, and Jorge Alberto~Olarte,
  \emph{Generalized permutahedra and positive flag dressians}, Int. Math. Res.
  Not. IMRN \textbf{2023} (2023), no.~19, 16748--16777.

\bibitem[JP18]{MR3829288}
Relinde Jurrius and Ruud Pellikaan, \emph{Defining the {$q$}-analogue of a
  matroid}, Electron. J. Combin. \textbf{25} (2018), no.~3, Paper No. 3.2, 32.
  \MR{3829288}

\bibitem[Mur10]{murota2010matrices}
Kazuo Murota, \emph{Matrices and matroids for systems analysis}, Springer,
  Berlin, 2010.

\bibitem[Wan26]{wang26modular}
Jidong Wang, \emph{Tropical projective duality of modular matroids, in
  progress}, 2026.

\bibitem[YY07]{yu2006representing}
Josephine Yu and Debbie Yuster, \emph{Representing tropical linear spaces by
  circuits}, Proc. FPSAC, 2007.

\end{thebibliography}


\end{document}
